\section{Ограничения применимости и обработка особых случаев}

Корректность расчётов определяется
областью применимости IF97 и устойчивостью численных процедур.
Поэтому в ядре выполняются
валидация входных параметров и явная обработка ошибок.

К типовым ограничениям относятся:
\begin{itemize}
\item допустимые диапазоны давления и температуры в соответствии со стандартом IAPWS-IF97~\cite{iapws_if97_2012};
\item требование $\rho > 0$ для маршрутов, использующих плотность;
\item требование $x \in [0; 1]$ для расчётов в двухфазной области;
\item ограничение области метастабильного пара в расширении региона~2 по давлению.
\end{itemize}

К особым случаям относятся вычисления на фазовых границах, в околокритической области и при переходах между регионами.
В этих ситуациях важны:
\begin{itemize}
\item контроль смены региона при итерационном процессе;
\item обработка случаев плохой обусловленности (например, вблизи критической точки);
\item ограничение числа итераций и явная сигнализация о несходимости;
\item единый формат ошибок валидации и вычислений, пригодный для интерфейсного и сервисного использования.
\end{itemize}

Перечисленные меры обеспечивают
устойчивую работу модели в инженерных сценариях
и воспроизводимость результатов в \texttt{if97\_core},
FLTK,
в оболочке на базе Yew, WebAssembly и Tauri
и в gRPC-сервисе.
